\documentclass[10pt,a4paper]{amsart}
\usepackage[margin=19mm]{geometry}
\usepackage{amsmath,amssymb,mathtools}
\usepackage[scr=rsfs,cal=euler]{mathalfa}
\usepackage{xfrac}
\usepackage[unicode,hidelinks,bookmarks=false]{hyperref}
\newcommand{\ie}{i.e.\ }
\newcommand{\cf}{cf.\ }
\newcommand{\resp}{respectively\ }
\newcommand{\Subsection}{\subsection}
\providecommand{\nobreakdash}{\nobreak\hbox{-}}
\setlength{\emergencystretch}{2em}
\multlinegap0pt
\usepackage{amssymb}
\usepackage{enumerate}
\newtheorem{lemma}{Lemma}
\title{Local-global principles for the existence of Levi factors}
\author{David Harbater, Julia Hartmann,, George McNinch}
\date{Source: arXiv:2603.26501v1 (CC BY 4.0)}
\begin{document}
\maketitle

\noindent\emph{Source and licence.} Local-global principles for the existence of Levi factors. Version arXiv:2603.26501v1 (CC BY 4.0), distributed by arXiv under the Creative Commons Attribution 4.0 International licence, http://creativecommons.org/licenses/by/4.0/, as recorded in the arXiv licence metadata for this exact version and shown on the abstract page https://arxiv.org/abs/2603.26501v1. Full licence notice: \texttt{licenses/arxiv-2603.26501-CC-BY-4.0.txt}.\par\smallskip

\noindent\textit{Editorial scope.} Counted unit: the article's lemma trivial (source0.tex lines 253--259). The article's complete original proof of it is delivered with it (source0.tex lines 260--263, one \texttt{proof} environment of 4 lines). No mathematics is written, repaired or re-derived: the statement and the proof are the article's own lines, in the article's own notation, and no retained line is substituted or re-flowed.  No further source-local context is needed: every cross-reference of every retained line resolves inside the two delivered spans.  Nothing was omitted from any retained span, so the package carries no omission record.  No retained line contains a citation, so the wrapper carries no bibliography and no bibliography entry.  No retained line contains a figure environment, an image-inclusion command or drawing code, so no image is delivered and the package declares no asset dependency.  Editorial formatting only, in addition: an AMS \texttt{amsart} preamble that restates the article's own declarations and macros used by the retained lines, namely the environments \texttt{lemma} on separate counters, together with the standard packages those lines need.  The retained environments print in this document as Lemma 1 in order of appearance, whereas the article prints the same items with the numbering of its own full document.  The complete native source of the article as distributed in the arXiv e-print for this exact version is bundled unchanged as \texttt{sources/197286/source0.tex}.
\begin{lemma}
Let $R$ be a complete local domain of characteristic $p>0$ with fraction field~$L$, and let $f\in L$. 
\begin{enumerate}
\item[(a)]\label{trivial} If $f$ is in the maximal ideal of $R$, the field extension of $L$ given by $Y^p-Y-f=0$ is trivial. 
\item[(b)]\label{nontrivial} If $f$ has valuation equal to ~$-1$, the field extension of $L$ given by $Y^p-Y-f=0$ is nontrivial. 
\end{enumerate}
\end{lemma}

\begin{proof}
If $f$ is in the maximal ideal of $R$, then $u:=-\sum\limits_{i=0}^\infty f^{p^i}$ is a well defined element of the complete ring $R$. This element satisfies $u^p-u=f$, and so the extension is trivial, proving the first assertion. 
For the second part, let $u$ be a nonzero element of $L$. If $u$ has nonnegative valuation, then so does $u^p-u$. On the other hand, if $u$ has valuation $m<0$, then $u^p-u$ has valuation $mp$. In summary, no element of the form $u^p-u$ has valuation $-1$, and so $Y^p-Y-f=0$ has no solution in $L$.
\end{proof}

\end{document}
